\section{Introduction}
\label{sec:intro}

Prediction sets are useful because they replace a single fallible prediction with a short list
that is likely to contain the true label, which a human can then verify. Consider a botanist
who photographs a specimen and receives a handful of candidate species: identifying the plant
from that list is easy, whereas confirming or rejecting one unqualified guess is not.
Conformal prediction turns any trained classifier into such a procedure, with a finite-sample
coverage guarantee and no assumptions on the data distribution
\cite{vovk2005algorithmic,angelopoulos2023gentle}.

The guarantee that comes for free is marginal: averaged over all inputs, the set contains the
true label with probability at least $1-\alpha$. Since marginal coverage is a prevalence-weighted
average of the class-conditional coverages, it is dominated by the frequent classes, and a
single shared threshold systematically under-covers the hard and rare ones. In biodiversity
monitoring or disease diagnosis the rare classes are often the ones that matter most, so the
species a botanist most needs to catch is the one whose prediction set most often omits it.
Recovering per-class control requires calibrating a separate threshold for each class, which is
what classwise conformal prediction \cite{vovk2005algorithmic,sadinle2019least} and its
successors \cite{ding2023class,shi2024rank,ding2026conformal} do.

Those methods all require labelled calibration examples of the class whose threshold they are
setting. In long-tailed deployments most rare classes are consumed by model training, which,
as \citet{ding2026conformal} note, leaves them with ``few or zero holdout examples to use for
calibrating uncertainty quantification methods.'' The zero case is common: on Pl@ntNet-300K
\cite{garcin2021plantnet} the median class has three held-out examples and on
iNaturalist-2018 \cite{vanhorn2018inaturalist} it has two. We measured what existing methods
emit for a class with no calibration examples, using the authors' released implementations,
and summarise the behaviours in \cref{tab:methods}. \textsc{Classwise} CP and \textsc{Rc3p}
return an infinite threshold, so the prediction set becomes the entire label space;
\textsc{Clustered} CP assigns the class to a null cluster and falls back to one shared
marginal threshold; and the kernel-weighted method of \citet{ding2026conformal}, which is the
closest approach to ours in structure, embeds each class using its own calibration scores and
therefore collapses every zero-example class onto one point of the embedding, handing them all
the same threshold. None of these behaviours is unreasonable, because with no labelled example
of a class there is nothing in the calibration set to condition on.

\paragraph{Objective.} Something about a class is nevertheless known without any of its
labels. A trained classifier already places every class somewhere in its output geometry: the
row of the classifier head that scores class $y$ has a norm and a bias, and it sits at some
set of distances from the rows of competing classes. Crowded rows correspond to classes that
are easily confused and therefore hard to cover. These quantities are available for a class
with zero calibration examples, since they are determined by the model's parameters and do
not involve any sample. Our goal is to use them to assign such a class a threshold that is specific to it,
while keeping the marginal coverage guarantee that split conformal prediction provides.

\paragraph{Our approach.} Predicted Class Correction (PCC) proceeds in three steps. On the
classes that do have calibration data we observe the offset $\delta_y$ between the per-class
conformal quantile and the marginal one, and fit a ridge regression of $\delta_y$ on the
descriptors $\varphi(y)$. We then apply the fitted regression to every class, including those
with no data, and shrink its output by a factor $\lambda$ chosen on the labelled classes only.
Finally we refit one scalar offset so that marginal coverage returns to its target. The
shrinkage step does substantial work here, since worst-class coverage is governed by the
largest error in the predicted offset, which its average error does not control, so applying
the raw prediction at full strength is harmful (\cref{tab:ablation}).

\paragraph{Summary of findings.} We evaluate on ImageNet \cite{russakovsky2015imagenet}, two
long-tailed datasets \cite{garcin2021plantnet,vanhorn2018inaturalist}, three backbones, four
score functions and $15$ corruption types \cite{hendrycks2019benchmarking}, for a total of
$228$ configurations. (i) On ImageNet classes with zero calibration examples, PCC improves
worst-class coverage by $0.075$ at matched average set size, against $0.027$ for the strongest
of seven competitors, five of which change the statistic by exactly zero
(\cref{tab:competitors}). (ii) Few classes need to be labelled for this to work: holding out
$90\%$ of classes, which leaves $100$ labelled classes out of $1000$, still yields $0.052$
with a confidence interval excluding zero, and the structural floor is $p+2$ labelled classes
for $p$ descriptors (\cref{fig:heldout}). (iii) Below roughly $35$ evaluation examples per
class, an oracle that observes the test labels cannot improve worst-class coverage either, so
what disappears in that regime is the room to improve and not the ability to take it
(\cref{fig:evaldepth}). This one axis accounts for all four settings in which the method
initially appeared to fail, and restricting evaluation to classes above the threshold recovers
a significant improvement on Pl@ntNet-300K. We arrived at it after discarding two explanations
of our own, both of which \cref{sec:discussion} records together with the measurements that
ruled them out.
